\chapter{Background and Theory}

\section{Zero-Knowledge Proofs}

Zero-knowledge proofs let a prover convince a verifier that a statement is true
without revealing secret witness data. Three properties are central in this
project:

\begin{itemize}
\item \textbf{Completeness:} honest executions produce proofs that verify.
\item \textbf{Soundness:} false claims cannot be proven except with negligible
probability.
\item \textbf{Zero-knowledge:} proofs leak no useful information about secret
inputs beyond validity.
\end{itemize}

In zkVM settings, these properties are applied to program execution rather than
to a standalone cryptographic statement: the prover shows that a fixed guest
binary ran correctly on specific inputs, and the verifier checks this without
replaying full execution. This direction continues a longer line of verifiable
computation work --- from circuit-based systems such as Pinocchio
\cite{parno2013pinocchio}, which compile fixed computations into arithmetic
constraints, to STARK-based proving for general execution traces
\cite{bensasson2018stark,risczeroStarkProtocol}, on which RISC Zero's
instantiation builds. Pinocchio
is cited here as foundational verifiable-computation prior art rather than as a
direct zkVM precursor in the RISC-V guest/host sense used by RISC Zero.

This distinction matters for the rest of the dissertation. The proof establishes
correct execution of code, while confidentiality and integrity claims still
depend on the underlying cryptographic construction and assumptions.

\section{RISC Zero Architecture}

RISC Zero uses a host/guest split \cite{risczeroDocs}.

\begin{itemize}
\item \textbf{Guest:} deterministic RISC-V program that performs cryptographic
computation and commits public outputs to the journal.
\item \textbf{Host:} prepares inputs, invokes proving, and verifies the receipt
against the expected image ID and journal output.
\end{itemize}

The proving pipeline is: compile guest code, execute guest with inputs, record
execution trace, generate receipt, then verify receipt validity and program
binding.

\subsection*{Competing zkVM Platforms and Selection Rationale}

Competing platforms exist, including SP1, Valida, and Plonky3. They differ in
proving backends, performance envelopes, and developer ergonomics.

For this project, RISC Zero was selected for practical reasons: a stable
host/guest contract (image ID plus journal), straightforward receipt
verification semantics, and tooling that supported repeated benchmark campaigns
within dissertation time constraints.

This choice is methodological, not universal. The report therefore treats any
platform-specific finding as scoped evidence rather than a claim about all
zkVMs.

\section{Cipher Design Preliminaries}

This section introduces the design characteristics of each cipher evaluated in
this project. Understanding their internal structure is necessary for
interpreting why they perform differently under zkVM proving.

\paragraph{AES (Advanced Encryption Standard).}
AES is a substitution-permutation network standardised by NIST in FIPS 197
\cite{nistfips197,daemen2002rijndael}. It operates on a 128-bit block and
supports key sizes of 128, 192, and 256 bits; this project uses AES-128. Each
round applies four operations: \texttt{SubBytes} (a nonlinear S-box
substitution), \texttt{ShiftRows} (byte permutation), \texttt{MixColumns}
(linear diffusion over GF($2^8$)), and \texttt{AddRoundKey} (XOR with the round
key). AES-128 has 10 rounds. The S-box and MixColumns layers are the main
engineering targets for implementation optimization; different representations
(lookup tables, bitslicing, combinational logic) give different trade-offs
between code size, data movement, and susceptibility to timing side channels
\cite{canright2005compactsbox,boyar2010logicmin}.

\paragraph{Salsa20.}
Salsa20 is an ARX (add-rotate-XOR) stream cipher designed by Bernstein
\cite{bernstein2008salsa20}. It generates keystream by applying a sequence of
quarter-round operations to a 512-bit state derived from a 256-bit key, a
64-bit nonce, and a 64-bit counter. Its design avoids S-boxes entirely:
nonlinearity comes from the interaction of addition (mod $2^{32}$) and rotation.
This makes the software path very regular and branch-free, which is favorable on
conventional CPUs but still accumulates significant trace cost in a zkVM because
of the number of state update steps.

\paragraph{LowMC.}
LowMC is a block cipher designed specifically to minimize the number of AND
gates (multiplicative depth and count) in arithmetic or Boolean circuit
representations \cite{albrecht2015ciphers,jaques2019grover}. Each round applies
a partial S-box layer (only a fraction of the state passes through nonlinear
substitution), a full-state linear layer (a random binary matrix multiply), and a
round-key XOR. The reduced nonlinearity was intended to make LowMC cheaper to
evaluate in MPC and ZK proof systems where AND gates are the dominant cost.
However, in a software zkVM, the large matrix multiplication in every round
becomes the dominant cost driver, as the matrices must be applied in guest code
and every operation contributes to the execution trace.

\section{Symmetric Cryptography in Scope}

This dissertation benchmarks three symmetric-cipher families in the same zkVM
measurement style:

\begin{itemize}
\item \textbf{AES-128,} a standardized block-cipher baseline
\cite{nistfips197,daemen2002rijndael};
\item \textbf{Salsa20,} an ARX stream-cipher baseline
\cite{bernstein2008salsa20};
\item \textbf{LowMC,} a design motivated by low nonlinear complexity in
MPC/ZK-oriented settings \cite{albrecht2015ciphers,jaques2019grover}.
\end{itemize}

The literature positions these choices differently, and that contrast motivates
the benchmark design:

\begin{itemize}
\item \textbf{AES:} not usually framed as ``ZK-specialized'', but its
standardization and mature implementation literature can make it competitive in
VM-oriented execution \cite{nistfips197,daemen2002rijndael}.
\item \textbf{Salsa20:} ARX operations are software-friendly, but long keystream
paths can still accumulate trace cost in zkVM contexts.
\item \textbf{LowMC:} favorable circuit-level nonlinearity profile, yet matrix
and bit-level software paths may impose substantial VM overhead.
\end{itemize}

For the final authenticated path, AES is used in CTR mode under the usual
nonce/IV-uniqueness requirement, following NIST Special Publication 800-38A's
CTR-mode specification \cite{nistsp80038a}.

Implementation-level AES literature is also relevant because S-box and
linear-layer engineering choices can materially change software behavior, even
when the algorithm is fixed \cite{canright2005compactsbox,boyar2010logicmin}.

\section{Related Work Positioning}

Prior ZK-focused primitive work is largely circuit-centric. LowMC-oriented
studies emphasize reduced nonlinear gate pressure
\cite{albrecht2015ciphers,jaques2019grover}, while other ZK-oriented primitives
are commonly evaluated by arithmetic-constraint efficiency rather than VM
execution behavior.

That body of work is valuable, but it leaves a gap for this dissertation's
question: how those design intuitions survive contact with an end-to-end zkVM
software stack. This report addresses that gap by using receipt-backed,
artifact-traceable measurements in one fixed execution model.

\section{Critical Comparison and Design Implications}

The related work does not point to a single universally best primitive because
each source optimizes for a different cost model. Circuit-oriented LowMC work
optimizes nonlinear constraint pressure; AES implementation work optimizes
software data paths and table or logic representations; zkVM documentation
defines a VM execution contract where proving follows the generated instruction
trace rather than a handwritten primitive circuit.

This creates the main literature gap used by the rest of the dissertation:
constraint-level efficiency does not automatically imply zkVM efficiency. A
primitive can have a favorable arithmetic or Boolean-circuit profile and still
lose if the guest implementation expands into many memory operations, branches,
or repeated bit transformations. Conversely, a conventional primitive such as AES
can be competitive if its software path is compact and stable in the VM.

\begin{table}[t]
\centering
\scriptsize
\begin{tabular}{p{0.16\linewidth}p{0.22\linewidth}p{0.22\linewidth}p{0.18\linewidth}}
\toprule
Evidence base & Main strength & Gap for this project & Design response \\
\midrule
LowMC/ZK-friendly cipher literature & Explains why low nonlinear complexity can
reduce circuit cost & Does not directly measure RISC Zero guest execution traces
& Benchmark LowMC rather than assume it wins \\[4pt]
AES implementation literature & Shows mature software engineering options for
S-box and linear-layer implementation & Usually targets CPU or constant-time
software goals, not proof cost & Keep AES as a serious zkVM baseline and test
backend variants \\[4pt]
zkVM platform documentation & Defines host/guest proof semantics and image-ID
binding & Does not decide which primitive is cheapest for a workload & Use a fixed
harness and artifact-backed comparisons \\[4pt]
Authenticated-encryption literature & Provides composition-level reasoning for
Encrypt-then-MAC & Does not bind transcripts to zkVM receipt objects by itself &
Combine EtM reasoning with receipt verification and journal commitments \\
\bottomrule
\end{tabular}
\caption{How background literature maps to design choices in this dissertation.}
\label{tab:literature-to-design}
\end{table}

The table also explains why later chapters use measurements as decision evidence
instead of treating the literature review as a direct ranking. The literature
defines expectations and security notions; the benchmark chapters test whether
those expectations survive in the implemented stack.

\section{Authenticated Encryption Framing}

Confidentiality alone is insufficient in adversarial settings, so the final
construction follows Encrypt-then-MAC: first encrypt with AES-CTR under NIST SP
800-38A, then authenticate public context and ciphertext with HMAC-SHA256 as
standardized by RFC 2104, FIPS 198-1, and FIPS 180-4
\cite{nistsp80038a,rfc2104,nistfips1981,nistfips1804}. In this implementation,
the emitted tag is truncated to 192 bits.

The choice of SHA-256 (used inside HMAC) is motivated by its origin in a
recent, externally vetted standard: FIPS 180-4, published by NIST
\cite{nistfips1804}. Using a standardised hash function means the security of
the MAC is grounded in a well-studied primitive rather than a custom or
ad-hoc construction. This is why HMAC-SHA256 was selected as the authentication
mechanism rather than, for example, a polynomial-based MAC or a cipher-based
construction.

RFC 2104 specifies the standard HMAC truncation rule as outputting at least 80
bits and recommends a minimum of half the hash output length (i.e.\ 128 bits
for SHA-256). This project uses a 192-bit truncation rather than the RFC minimum
of 128 bits, for the following reason: AES-128 offers a 128-bit key security
level, while the HMAC tag determines the forgery resistance. Truncating to 192
bits means the tag provides a \(2^{192}\) forgery work factor for a blind
attempt, which is tighter than the 128-bit AES key security. The tag length
is therefore the security-determining link in the chain rather than the AES key
length: an adversary who cannot forge the HMAC tag cannot tamper with the
authenticated transcript, regardless of the AES key size. Using 192 bits
explicitly acknowledges this and provides a larger margin than the RFC minimum
without incurring the full 256-bit output overhead.

This framing is supported by standard composition results for authenticated
encryption reasoning \cite{bellare2000authenticated}.

Within this project, the framing appears as two benchmarked mode families:
CTR-only (\texttt{aes-ctr-*}) and CTR+HMAC (\texttt{aes-ctr-hmac-*}). This
separation is implementation-oriented and supports direct overhead attribution in
later evaluation chapters.

The main notions used in the write-up are:

\begin{itemize}
\item \textbf{IND-CPA} for confidentiality against chosen-plaintext attacks
in Chapter 10's proof-output challenge game;
\item \textbf{HMAC PRF/MAC assumptions} for tag indistinguishability and
tamper-detection reasoning;
\item \textbf{zkVM proof properties} for execution binding and non-journaled
witness-data leakage assumptions.
\end{itemize}
